% DeXposure-FM architecture section (drop-in LaTeX).
% Requires in the preamble:
% \usepackage{tikz}
% \usetikzlibrary{arrows.meta,positioning,fit}
%
% This section aligns notation with FormalisingCreditExposures-clean.tex:
% exposure definition: Eq.~\eqref{eq:credit-exposure-definition}
% node weight: Eq.~\eqref{eq:node-weight}
% value flow: Eq.~\eqref{eq:value-flow}
% edge weight: Eq.~\eqref{eq:edge-weight}

\subsection{DeXposure-FM Architecture}
\label{sec:dexposure-fm-architecture}

DeXposure-FM is trained for multi-step forecasting of edge existence, edge weights, and node TVL changes. %; and (2) forward-looking financial stability assessment by applying risk metrics and stress tests to predicted exposure graphs, and macroprudential tools in Section~\ref{sec:econ}. We provide a reference implementation of these tools (observed and forecast-then-measure) in \texttt{run\_macroprudential\_tools.py}. Unlike traditional financial networks where data completeness is a concern, blockchain-based DeFi networks provide fully observable on-chain transaction data, eliminating the need for edge imputation tasks common in opaque interbank settings.

\subsubsection{Inputs and Output Settings}

\begin{comment}
Let $t\in\mathcal{T}$ denote discrete observation times and let $\tau = t_2-t_1$ be a discrete interval between consecutive times.
For each interval $\tau$ we construct a weighted directed graph snapshot
\begin{equation}
G_{\tau}=(\mathcal{P}_{\tau},\mathcal{E}_{\tau}),
\end{equation}
where $\mathcal{P}_{\tau}$ is the set of protocols and $\mathcal{E}_{\tau}$ is the set of directed exposure links.
Node weights are given by $w_{\tau}(p)$ (the protocol's end-of-interval stock) as in Eq.~\eqref{eq:node-weight}, and edge weights are given by $w_{\tau}(e_{pq})$ (the change in credit exposure from $p$ to issuing protocol $q$ over $\tau$) as in Eq.~\eqref{eq:edge-weight}, where token-level value flows $F^{\sigma}_{pq,\tau}$ are defined by Eq.~\eqref{eq:value-flow}. Credit exposure exists when $X_p(t)\cap X_G(q)\neq\varnothing$ (Eq.~\eqref{eq:credit-exposure-definition}).
\end{comment}

%For each protocol $p$ we form the following inputs:

\fh{what are models' IO?}

The inputs for the DeXposre-FM is as follow.s
\begin{itemize}
  \item \emph{graph snapshot:} $G_{\tau}$ with node features including $w_{\tau}(p)$ and edge features including $w_{\tau}(e_{pq})$ (and optional attributes such as token class or sector pair).
  \item \emph{tabular node features:} $\mathbf{x}^{\mathrm{tab}}_{p}$ contains protocol descriptors including:
  \begin{itemize}
    \item Log-scaled TVL: $\log(1 + w_{\tau}(p))$
    \item Number of token types held
    \item Maximum token share (concentration measure)
    \item Token composition entropy
    \item Protocol category (one-hot encoded, $\sim$15 classes)
    \item Previous week TVL change: $\log(1 + w_{\tau}(p)) - \log(1 + w_{\tau-1}(p))$ (temporal momentum)
    \item Normalized in-degree: incoming edge count / max in-degree
    \item Normalized out-degree: outgoing edge count / max out-degree
  \end{itemize}
\end{itemize}

\subsubsection{Graph-Tabular Encoder and Weight Prior}
We employ GraphPFN \cite{eremeev2025turningtabularfoundationmodels} %, a pre-trained graph-tabular foundation model based on Transformer architecture, 
as a graph-tabular encoder in our model. %GraphPFN jointly processes graph structure and tabular node features through multi-head self-attention with graph structure injection.
A LiMiX transformer %\fh{\cite{xx}} %within GraphPFN 
produces %\fh{per-feature embeddings} 
for each node. For link prediction, we extract node representations by mean-pooling %\fh{\cite{xx}} 
over the \textit{feature column embeddings} %\fh{\cite{xx}} 
rather than using the label column embedding %\fh{\cite{xx}} 
(which is designed for node classification):
\begin{equation}
\mathbf{h}_{p,\tau} = \frac{1}{F}\sum_{f=1}^{F} \mathbf{e}_{p,\tau}^{(f)},
\end{equation}
where $\mathbf{e}_{p,\tau}^{(f)}\in\mathbb{R}^d$ is the embedding for feature $f$ of protocol $p$, and $F$ is the number of tabular features. This aggregation captures the joint representation of all node attributes, which is more suitable for edge-level tasks than the label-specific embedding.

\begin{comment}
    
The encoder can be operated in two modes:
\begin{itemize}
  \item \textbf{Frozen probing:} Encoder weights are frozen; only task heads are trained. This measures the transferability of pre-trained representations.
  \item \textbf{Fine-tuned:} End-to-end training with differential learning rates---lower rate ($\eta/10$) for the pre-trained backbone and higher rate ($\eta$) for task-specific heads.
\end{itemize}
\end{comment}

\subsubsection{Task Heads}
\textbf{(1) Multi-step forecasting.}
For horizon $h\in\{1,4,8,12\}$ weeks, we predict future edge existence and weights. The link scorer uses a Hadamard-product MLP architecture:
\begin{align}
\mathbf{f}_{pq} &= [\mathbf{h}_{p,\tau}; \mathbf{h}_{q,\tau}; \mathbf{h}_{p,\tau}\odot\mathbf{h}_{q,\tau}; |\mathbf{h}_{p,\tau}-\mathbf{h}_{q,\tau}|],\\
(\widehat{y}^{\mathrm{exist}}_{pq}, \widehat{w}_{\tau+h}(e_{pq})) &= \mathrm{MLP}_{\mathrm{link}}(\mathbf{f}_{pq}),
\end{align}
where $\widehat{y}^{\mathrm{exist}}_{pq}\in(0,1)$ is the predicted probability of edge existence, and $\widehat{w}_{\tau+h}(e_{pq})$ is the predicted edge weight (log-scaled).

For node-level prediction, we forecast the log-change in protocol TVL:
\begin{equation}
\widehat{\Delta}_{p,\tau+h} = \mathrm{MLP}_{\mathrm{node}}(\mathbf{h}_{p,\tau}),
\end{equation}
where $\widehat{\Delta}_{p,\tau+h} = \log(1+w_{\tau+h}(p)) - \log(1+w_{\tau}(p))$ for protocols $p\in\mathcal{P}_{\tau}\cap\mathcal{P}_{\tau+h}$.

Network-level statistics (concentration, density, sector connectivity) are obtained by applying deterministic functionals to the predicted graph $\widehat{G}_{\tau+h}$.

\textbf{(2) Shock analysis.}
We evaluate model performance during extreme market events---specifically the Terra/Luna collapse (2022-05-09) and FTX collapse (2022-11-07). For each event, we examine a window around the shock date to assess:
\begin{itemize}
  \item Network structure changes (TVL, edge count, Gini, HHI)
  \item Forward-looking risk trajectories on predicted graphs
  \item Early warning signal detection capability
\end{itemize}


\begin{comment}
    
\begin{figure}[t]
\centering
\begin{tikzpicture}[
  font=\footnotesize,
  scale=0.88, transform shape,
  box/.style={draw, rounded corners, align=center, minimum width=3.0cm, minimum height=0.9cm},
  sbox/.style={draw, rounded corners, align=center, minimum width=2.75cm, minimum height=0.82cm},
  arrow/.style={-Latex, thick},
  node distance=0.5cm and 0.85cm
]
  % Inputs
  \node[box] (tabin) {Tabular node features\\ $\mathbf{x}^{\mathrm{tab}}_{p}$};
  \node[box, below=of tabin] (graphin) {Exposure graph snapshot\\ $G_{\tau}=(\mathcal{P}_{\tau},\mathcal{E}_{\tau})$\\ Node: $w_{\tau}(p)$, Edge: $w_{\tau}(e_{pq})$};

  % Encoder
  \node[sbox, right=1.5cm of graphin] (egraph) {Pre-trained\\GraphPFN encoder\\ $E_{\mathrm{GraphPFN}}$};

  % Node embeddings
  \node[sbox, right=1.3cm of egraph] (embed) {Node embeddings\\ $\mathbf{h}_{p,\tau}$};

  % Heads
  \node[sbox, right=1.2cm of embed] (forecast) {Forecasting head\\ $\widehat{w}_{\tau+h}(e_{pq}),\,\widehat{\Delta}_{p,\tau+h}$};
  \node[sbox, above=0.55cm of forecast] (shock) {Shock analysis\\ Terra/FTX events};
  % Arrows
  \draw[arrow] (tabin.east) -- ++(0.55,0) |- (egraph.west);
  \draw[arrow] (graphin) -- (egraph);

  \draw[arrow] (egraph) -- (embed);
  \draw[arrow] (embed) -- (forecast);

  \draw[arrow] (embed.east) -- ++(0.3,0) |- (shock.west);

  % Grouping boxes
  \node[draw, dashed, rounded corners, fit=(tabin)(graphin), inner sep=6pt, label={[align=center]above:Inputs}] {};
\end{tikzpicture}
\caption{DeXposure-FM architecture aligned with the DeXposure exposure-graph construction. The model uses a pre-trained GraphPFN encoder to jointly process graph structure and tabular node features, producing node embeddings $\mathbf{h}_{p,\tau}$ that support multi-step forecasting and shock analysis tasks.}
\label{fig:dexposure-fm-architecture}
\end{figure}

\end{comment}
